%%%%%%%%%%  scan idx 293  |  folio 280  %%%%%%%%%%
\item Le
groupe \(K_\bullet({}_A X)_{\naif}\) n'a pratiquement d'intérêt que si \(A\) est cohérent (autrement il risque d'y avoir fort peu de faisceaux cohérents \ldots).
\item La considération du \(k\) de catégories triangulées «~trop grosses~» n'offre pas d'intérêt. Ainsi, de (1.2 b)) on déduit facilement les relations
\[
 k(D({}_A X))=k(D^-({}_A X))=k(D^+({}_A X))=k(D^b({}_A X))=k(D^-({}_A X)_{\coh})=0
\]
(pour la dernière égalité, observer que \(L\in\Ob D^-({}_A X)_{\coh}\) implique \(\coprod_{i\geq0}L[2i]\in\Ob D^-({}_A X)_{\coh}\)).
\end{enumerate}

\textbf{2.4. Cas d'un faisceau d'anneaux cohérent.} Soit \((X,A)\) un topos annelé, où \(A\) est supposé \emph{cohérent} en tant que \(A\)-Module à gauche (I 3.1 et 3.6). Tout \(A\)-Module à gauche cohérent est un complexe pseudo-cohérent (I 3.5), on a donc un foncteur pleinement fidèle
\begin{equation}\tag{2.4.1}
 \Coh({}_A X)\hookrightarrow D^b({}_A X)_{\coh}.
\end{equation}
Or (I 3.4) \(\Coh({}_A X)\) est une sous-catégorie épaisse de \(\Mod({}_A X)\), et, en vertu de (I 3.5), \(D^b({}_A X)_{\coh}\) est la sous-catégorie pleine de \(D^b({}_A X)\) formée des objets \(L\) tels que \(H^i(L)\in\Ob\Coh({}_A X)\) pour tout \(i\). Par suite, il résulte de (1.6) que \emph{l'homomorphisme}
\begin{equation}\tag{2.4.2}
 K_\bullet({}_A X)_{\naif}\longrightarrow K_\bullet({}_A X)
\end{equation}
\emph{défini par} (2.4.1) \emph{est un isomorphisme}.

\textbf{2.5. L'homomorphisme \(K^{\bullet}\to K_\bullet\).} Soit \((X,A)\) un topos annelé. Le foncteur d'inclusion \(D({}_A X)_{\uparf}\hookrightarrow D^b({}_A X)_{\coh}\) définit un homomorphisme
\begin{equation}\tag{2.5.1}
 K^{\bullet}({}_A X)\longrightarrow K_\bullet({}_A X).
\end{equation}
Celui-ci n'est en général ni injectif ni surjectif. Cependant, faisons les hypothèses suivantes :
\begin{enumerate}[label=(\roman*)]
